\section{System Architecture}

The service runs in Node.js, connects to MongoDB through the official Node driver, and uses Groq-hosted open models for generation. A request passes through authentication, connection retrieval, schema profiling, ambiguity gating, optional clarification, generation, validation, and execution with a 30\,s server-side time limit. In the default \texttt{hybrid} mode the model may itself return clarification questions; these are accepted only if they offer 2--4 fields that exist in the schema.

\begin{figure*}[t]
\centering
\begin{tikzpicture}[font=\footnotesize,
  box/.style={draw, rounded corners=2pt, minimum height=7mm, minimum width=18mm, align=center, fill=gray!8},
  store/.style={draw, minimum height=7mm, minimum width=18mm, align=center, fill=blue!6},
  arr/.style={-{Stealth[length=2mm]}, thick}, node distance=6mm and 4mm]
\node[box] (client) {Client\\(React)};
\node[box, right=of client] (api) {API gateway\\auth, scope};
\node[box, right=of api] (prof) {Schema\\profiler};
\node[box, right=of prof, fill=orange!15] (gate) {Ambiguity\\gate};
\node[box, right=of gate] (llm) {LLM\\generation};
\node[box, right=of llm, fill=green!10] (val) {Safety\\validator};
\node[box, right=of val] (exec) {Executor};
\node[box, below=of gate, fill=orange!15] (card) {Clarification\\card};
\node[box, below=of llm] (enr) {Allowlist\\+ enrich};
\node[store, below=of api] (vault) {Vault + pool};
\node[store, above=of llm] (mem) {Few-shot\\memory};
\node[box, below=of val, fill=red!8] (hitl) {Approval token\\(write staging)};
\draw[arr] (client)--(api); \draw[arr] (api)--(prof); \draw[arr] (prof)--(gate);
\draw[arr] (gate)--node[above]{clear}(llm);
\draw[arr] (llm)--(val); \draw[arr] (val)--node[above]{safe}(exec);
\draw[arr] (gate)--node[left]{ambiguous}(card);
\draw[arr] (card)--(enr); \draw[arr] (enr)--(llm);
\draw[arr] (mem)--(llm);
\draw[arr] (val)--node[right]{mutating}(hitl);
\draw[arr, dashed] (vault.north east)--(prof.south west);
\end{tikzpicture}
\caption{AMB-MQL request flow. The vault and pool supply database handles to the profiler and executor (dashed).}
\label{fig:arch}
\end{figure*}

\subsection{Schema profiling and scope}
For each in-scope collection the profiler draws $N{=}5$ documents with \texttt{\$sample}, records field names, inferred BSON types and one sample value, flattens nested sub-documents into dot paths (depth at most three), and exposes sub-fields of arrays of objects as \texttt{field[].sub}. Collections whose names contain ``embed'' are skipped. The result is cached for five minutes, keyed by connection, database and sample size, and rendered as a compact schema string for the prompt. Because five documents can miss sparse fields, a field present in a fraction $p$ of documents is missed with probability $(1-p)^5$ under sampling with replacement (about $0.77$ for $p{=}0.05$). A user may restrict a request to a subset of collections; this scope bounds both the schema shown to the model and the collections the validator permits.

\subsection{Generation}
The system prompt contains the current date, the schema string, up to three retrieved question--pipeline examples, the output format and aggregation rules (for example, group before project, unwind arrays before grouping). Examples are stored per tenant in the central database and retrieved with an Atlas Search text query, falling back to keyword regular-expression matching. A new example is stored after an execution returns a non-empty result; no further verification is applied (Section~\ref{sec:limitations}).

\subsection{Validation, execution and write staging}
Every generated pipeline passes a recursive key-scan validator. It rejects code-execution operators (\texttt{\$function}, \texttt{\$accumulator}, \texttt{\$where}) anywhere in the pipeline, enforces a 20-stage cap, requires each top-level stage to belong to an 18-stage analytical allowlist, resolves collections referenced through \texttt{\$lookup}, \texttt{\$unionWith}, \texttt{\$graphLookup} and \texttt{\$facet} and rejects those outside the scope, and appends \texttt{\$limit: 1000} when no limit is present. Pipelines with mutating stages are not executed. They are staged behind an approval token: an HS256-signed JWT valid for ten minutes that binds the pipeline, collection, scope and session, carries a unique identifier consumed once through a unique index, and is re-validated when submitted. The token is returned to the requesting session, so it provides confirmation rather than separation of duties. Results are normalised (Decimal128, ObjectId, dates) before being returned.

\subsection{Multi-tenant infrastructure}
Connection strings are encrypted with AES-256 in CBC mode~\cite{nist80038a} under a fresh random 16-byte IV; the 256-bit key is read from the environment and the service refuses to start without it. A connection is identified by the SHA-256 digest of its URI. Sessions use an httpOnly cookie holding a seven-day JWT, and requests are rate-limited per IP. Database handles are pooled per connection (at most 50 clients, five-minute idle timeout, least-recently-used eviction, a ping before reuse). CBC provides confidentiality without an integrity tag and the digest is unsalted; Section~\ref{sec:limitations} lists both.
